\documentclass[10pt,color,leqno]{beamer}
\usepackage[french]{babel}
\usepackage[utf8]{inputenc}
\usepackage{t1enc,times}
\usepackage{amssymb, amsmath}
\usepackage{appendixnumberbeamer}

\usefonttheme{serif}
\usetheme{JLTree}

\setbeamerfont{section in head/foot}{family=\sffamily}
\setbeamerfont{subsection in head/foot}{family=\sffamily}
\setbeamerfont{title in head/foot}{family=\sffamily}
\setbeamerfont{date in head/foot}{family=\sffamily}
\setbeamerfont{author in head/foot}{family=\sffamily}
\setbeamerfont{institute in head/foot}{family=\sffamily}
\setbeamercovered{invisible}


\theoremstyle{plain}
\newtheorem{theo}{Théorème}[section]


\def\ca{{\mathcal A}}
\def\cc{{\mathcal C}}
\def\cd{{\mathcal D}}
\def\ce{{\mathcal E}}
\def\cf{{\mathcal F}}
\def\cg{{\mathcal G}}
\def\ch{{\mathcal H}}
\def\ci{{\mathcal I}}
\def\cl{{\mathcal L}}
\def\cm{{\mathcal M}}
\def\cn{{\mathcal N}}
\def\ct{{\mathcal T}}
\def\cs{{\mathcal S}}
\def\cx{{\mathcal X}}
\def\cu{{\mathcal U}}
\def\cw{{\mathcal W}}
\def\C{{\mathbb C}}
\def\D{{\mathbb D}}
\def\E{{\mathbb E}}
\def\H{{\mathbb H}}
\def\I{{\mathbb I}}
\def\L{{\mathbb L}}
\def\N{{\mathbb N}}
\def\P{{\mathbb P}}
\def\Q{{\mathbb Q}}
\def\R{{\mathbb R}}
\def\Z{{\mathbb Z}}
\def\s{\star}
\def\l{\ell}

\newcommand{\ind}{{\bf 1}}
\newcommand{\expp}[1]{\mathop {\mathrm{e}^{ #1}}}
\newcommand{\Supp}{{\rm Supp}\;}
\def\esssup{\mathop{\mathrm{ess sup}}\nolimits}
\def\Var{\mathop{\mathrm{Var}}}
\def\Cov{\mathop{\mathrm{Cov}}}

\def\abs#1{\left|#1\right|}
\def\inv#1{\frac{1}{#1}}
\def\ind#1{{\bf 1}_{\left\{#1\right\}}}
\def\htau{\widehat{\tau}}
\def\htheta{\widehat{\theta}}

%% Beamer Layout %%%%%%%%%%%%%%%%%%%%%%%%%%%%%%%%%%
\useoutertheme[subsection=false,shadow]{miniframes}
\useinnertheme{default}
\usefonttheme{serif}

% -------------------------------------------------------------%
\def\red#1{{\color{red} #1}}
\def\blue#1{{\color{blue} #1}}
\setlength{\columnseprule}{1pt}
\definecolor{rose}{rgb}{0.92,0.10,0.12}
\definecolor{blue}{rgb}{0,0,1}

\setbeamertemplate{page number in head/foot}[totalframenumber]
\setbeamertemplate{theorem begin}
{%\setbeamercolor{block body}{bg=red}
\begin{\inserttheoremblockenv}
  {%
  \inserttheoremheadfont
  \inserttheoremname
  \inserttheoremnumber
  \ifx\inserttheoremaddition\empty\else\ (\inserttheoremaddition)\fi%
      % \inserttheorempunctuation
  }%
  }
  \setbeamertemplate{theorem end}
  {%
\end{\inserttheoremblockenv}}

\setbeamertemplate{bibliography item}[triangle]

\setbeamertemplate{frametitle continuation}{(\insertcontinuationcount)}
\setbeamertemplate{enumerate items}[ball]
\setbeamertemplate{itemize items}[triangle]

\newtheorem*{thm}{Theorem}
\newtheorem*{prop}{Proposition}
\newtheorem*{proposition}{Proposition}
\newtheorem*{cor}{Corollary}
\newtheorem{lem}{Lemme}
\newtheorem{lemmafr}{Lemme}
\newtheorem*{Algo}{Algorithm}

\definecolor{vert}{rgb}{0.12,0.90,0.}
\newenvironment{algo}[1][]{%
\setbeamercolor{block title}{fg=black,bg=green!40!black}
\setbeamercolor{block body}{parent=normal text,use=block title,bg=block title.bg!20!bg}
\begin{Algo}[#1]}{
\end{Algo}
\setbeamercolor{block title}{use=structure,fg=structure.fg,bg=structure.fg!20!bg}
\setbeamercolor{block body}{parent=normal text,use=block title,bg=block title.bg!50!bg}
}
\newenvironment{remark}{{\it Remark.} }{}

\def\s{\star}
%%%%%%%%%%%%%%%%%%%%%%%%%%%%%%%%%%%%%%%%%%%%%%%%%%

\title{Conception et implémentation d'un outil de couverture de produits dérivés en C++}
\author[J. Lelong]{Jérôme Lelong}
\date{2026-2027}



\begin{document}

\frame{\titlepage}

% --------------------------------------------------------SLIDE -
\begin{frame}
  \frametitle{Plan}
  \tableofcontents
\end{frame}
% --------------------------------------------------------------- %

\section{Objectifs}

% --------------------------------------------------------SLIDE -
\begin{frame}[allowframebreaks]
  \frametitle{Objectifs}

  \begin{itemize}
    \item Évaluer et couvrir des options path-dependent dans un modèle Black-Scholes multidimensionnel.
    \item Nous considérons des options avec des payoffs de la forme
    \begin{equation*}
      \varphi(S_{t_0}, S_{t_1}, \dots, S_{t_N})
    \end{equation*}
    où $0 = t_0 < t_1 < \cdots < t_N = T$ est une grille régulière de dates d'observation, c'est-à-dire $t_k = kT / N$ pour $k=0,\dots,N$.
    \item Quelques exemples
    \begin{itemize}
      \item Option d'achat: $N=1$ et $\varphi(S_0, S_1) = (S_1 - K)_+$.
      \item Option barrière discrète: $T = 1$, $N=12$ et $\varphi(S_0, S_1, \dots, S_{12}) = (S_{12} - K)_+ \ind{\forall 0 \le i \le 12,\; S_i \ge L}$.
    \end{itemize}
    \pagebreak
    \item Construire un portefeuille $V$ tel que
    \begin{align*}
      V_T &= \varphi(S_{t_0}, S_{t_1}, \dots, S_{t_N}) \\
      V_t &= H^0_t S^0_t + H_t \cdot S_t \\
      dV_t &= H^0_t dS^0_t + H_t \cdot dS_t
    \end{align*}
    où $H^0$ est la quantité d'actif sans risque $S^0_t = \expp{r t}$ et $H$ est la quantité d'actifs risqués.
    \item En temps continu, un tel $V$ existe grâce au théorème de représentation des martingales
    \begin{align*}
      V_t &= \expp{-r (T-t)} \E[ \varphi(S_{t_0}, S_{t_1}, \dots, S_{t_N}) | \cf_t] := u(t, S_t) \\
      H_t &= \partial_d u(t, S_t).
    \end{align*}
    \item Concevoir et implémenter un outil C++ pour construire une approximation discrète de $V$.
  \end{itemize}

\end{frame}
% --------------------------------------------------------------- %

% --------------------------------------------------------SLIDE -
\begin{frame}
  \frametitle{Les approximations numériques}

  \begin{itemize}
    \item Approcher les espérances
    \begin{itemize}
      \item Utiliser la moyenne empirique
      \begin{align*}
        \E[ \varphi(S_{t_0}, S_{t_1}, \dots, S_{t_N}) ] \approx \frac{1}{M} \sum_{m=1}^M \varphi(S^{(m)}_{t_0}, S^{(m)}_{t_1}, \dots, S^{(m)}_{t_N})
      \end{align*}
      \item Toujours fournir un intervalle de confiance
      \begin{align*}
      & \xi_M^2 = \left[\inv{M} \sum_{m=1}^M (\varphi(S_{t_0}^{(m)},  \dots, S_{t_N}^{(m)}))^2 - \left(\inv{M} \sum_{m=1}^M \varphi(S_{t_0}^{(m)},  \dots, S_{t_N}^{(m)}) \right)^2 \right] \\
      &\text{IC 95\%:} \left[ \inv{M} \sum_{m=1}^M \varphi(S_{t_0}^{(m)},  \dots, S_{t_N}^{(m)}) \pm \frac{1.96 \xi_M}{\sqrt{M}} \right].
    \end{align*}
    \end{itemize}
      \item Utiliser une méthode de différences finies pour approcher les dérivées partielles.
  \end{itemize}



\end{frame}
% --------------------------------------------------------------- %


\section{Le modèle Black-Scholes multidimensionnel}

% --------------------------------------------------------SLIDE -
\begin{frame}
  \frametitle{Le modèle Black-Scholes multidimensionnel}

  \begin{equation*}
    S_{t,d} = S_{0,d} \expp{(r- {(\sigma_d)}^2/2) t + \sigma_d B_{t,d}}, \quad d=1,\dots, D
  \end{equation*}
  où $\sigma_d$ sont des volatilités à valeurs réelles, $B_d$ sont des mouvements browniens à valeurs réelles pour tous $1 \le d \le D$ et $\frac{\Cov(B_t, B_t)}{t} = \Gamma$. \\

  Il existe $L$ tel que $LL' = \Gamma$. Alors,
  \begin{equation*}
    S_{t,d} = S_{0,d} \expp{(r- {(\sigma_d)}^2/2) t + \sigma_d L_d W_t}, \quad d=1,\dots, D
  \end{equation*}
  où $L_d$ désigne la $d^{\text{ème}}$-ligne de la matrice $L$ et $W = (W_1, \dots, W_D)$ est un mouvement brownien à valeurs dans $\R^D$. Les composantes de $W$ sont indépendantes.

\end{frame}
% --------------------------------------------------------------- %

% --------------------------------------------------------SLIDE -
\begin{frame}
  \frametitle{Simulation du modèle Black-Scholes}

  \begin{itemize}
    \item Notez que pour tous $t, u \ge 0$
      \begin{equation*}
        S_{t+u,d} = S^d_{t,d} \expp{(r- \sigma^2/2) u + \sigma (B_{t+u,d} - B_{t,d})} = S_t \tilde S_u
      \end{equation*}
    \item Considérez une grille régulière $t_k = \frac{kT}{H}$ pour $k=0,\dots,H$
      \begin{align*}
        S_{t_{i+1},d} & \stackrel{\text{Loi}}{=} S_{t_i,d} \expp{(r- {(\sigma_d)}^2/2) \frac{T}{H} + \sigma_d \sqrt{\frac{T}{H}} L_d G_{i+1}}, \quad d=1,\dots, D
      \end{align*}
      où $(G_i)_{i \ge 1}$ est iid $\cn(0, I)$.
    \item Soit $t_{i} \le t < t_{i+1}$, nous considérons le problème de simulation $(S_{t+u}, u \ge 0)$ sachant $(S_s, s \le t)$. Pour $d=1,\dots, D$,
      \begin{align*}
        S_{t_{i+1},d} & \stackrel{\text{Loi}}{=} S_{t,d} \expp{(r- {(\sigma_d)}^2/2) (t_{i+1} - t) + \sigma_d \sqrt{(t_{i+1} - t)}L_d G_{i+1}}                       \\
        S_{t_{k+1},d} & \stackrel{\text{Loi}}{=} S_{t_k,d} \expp{(r- {(\sigma_d)}^2/2) \frac{T}{H} + \sigma_d \sqrt{\frac{T}{H}} L_d G_{k+1}}, \quad k=i+1, \dots, D
      \end{align*}
  \end{itemize}

\end{frame}
% --------------------------------------------------------------- %

\section{Le portefeuille en temps discret}

% --------------------------------------------------------SLIDE -
\begin{frame}
  \frametitle{Le portefeuille en temps discret}

  \begin{itemize}
    \item Considérez une grille de temps régulière $0 = \tau_0 < \cdots < \tau_H = T$ avec un pas de temps $T/H$.
    \item Soit $\hat{p}_0 \approx u(0, S_0)$ et $\delta_i \approx \nabla u(\tau_i, S_{\tau_i})$.
    \item La valeur du portefeuille de couverture en temps discret est donnée par
      \begin{equation*}
        V_i = \delta_i \cdot S_{\tau_i} + \phi^0_i S^0_{\tau_i}
      \end{equation*}
      où $\phi^0_i$ est déterminé en utilisant l'équation d'auto-financement
      \begin{equation*}
        (V_{i+1} - V_i) = \delta_i \cdot (S_{\tau_{i+1}} - S_{\tau_{i}}) + \phi^0_i (S^0_{\tau_{i+1}} - S^0_{\tau_{i}})
      \end{equation*}
    \item Introduisons $\tilde V_i = \frac{V_i}{S^0_{\tau_i}}$. Alors,
      \begin{equation*}
        (\tilde V_{i+1} - \tilde V_{i}) = \delta_i \cdot (\tilde S_{\tau_{i+1}} - \tilde S_{\tau_i}).
      \end{equation*}
    \item L'erreur de couverture à la date d'échéance est donnée par $P\&L = \hat{p}_0 + \expp{r \tau_H} \tilde V_{H} - \text{payoff}$.
  \end{itemize}

\end{frame}
% --------------------------------------------------------------- %

\section{En pratique}

% --------------------------------------------------------SLIDE -
\begin{frame}
  \frametitle{En pratique}

  \begin{itemize}
    \item Projet à réaliser par groupe de 4.
    \item Il est interdit de chercher à faire réaliser le projet par une IA.
    \item Lire le sujet.
    \item Séance 1: réfléchir à l'architecture.
    \item Une architecture sera mise à dispos à partir de la séance 2.
  \end{itemize}

\end{frame}
% --------------------------------------------------------------- %

% --------------------------------------------------------SLIDE -
\begin{frame}
  \frametitle{Quelques bonnes pratiques}

  \begin{itemize}
    \item Écrivez des tests! Les tests font partie du code.
    \item Testez soigneusement chaque fonction.
    \item Documentez votre code.
    \item Ne cherchez pas à sur-optimiser votre code au détriment de sa lisibilité et maintenance.
    \item Utilisez l'option d'achat vanille
    \begin{itemize}
      \item pour vérifier votre implémentation Monte-Carlo,
      \item pour vérifier votre approximation par différences finies,
      \item pour vérifier votre équation de portefeuille.
    \end{itemize} 
  \end{itemize}

  

\end{frame}
% --------------------------------------------------------------- %
\end{document}
